\documentclass[a4paper,11pt]{article}
\usepackage[utf8]{inputenc}
\usepackage[T1]{fontenc}
\usepackage[english]{babel}
\usepackage{geometry}
\geometry{left=2cm,right=2cm,top=2cm,bottom=2cm}

\usepackage{lmodern}
\usepackage{xcolor}
\definecolor{titleblue}{RGB}{0,51,140}
\usepackage{hyperref}
\usepackage{titlesec}

\titleformat{\section}{\color{titleblue}\large\bfseries}{\thesection}{1em}{}

\begin{document}

\begin{center}
    {\LARGE\bfseries\color{titleblue} Research Statement}\\[6pt]
    {\large AI-Assisted Model-Driven Security Engineering for Digital Twins}\\[4pt]
    \textbf{Ismail AISSAOUI} $|$ State Engineer in AI
\end{center}

\bigskip

\section{Introduction and Motivation}
The design of complex cyber-physical digital twins requires the integration of security concerns from the earliest architectural stages. However, current security engineering remains largely manual and lacks synchronization between design-time models and runtime execution. This research statement proposes an **AI-assisted Model-Driven Engineering (MDE)** framework that closes the loop between design and execution, enabling automated security injection, synchronization, and simulation within industrial digital twins.

\section{Current Research: Relational Intelligence for Security}
While my recent work at Sonatrach focuses on physics-informed modeling, I have built a strong foundation in **Graph Neural Networks (GNNs)** for structural intelligence. My "Career Connect" and "Geo-RAG" projects involved modeling complex relational systems to extract semantic knowledge and detect structural anomalies. This experience is directly applicable to the challenge of mapping security standards (e.g., MITRE ATT\&CK) onto the relational models of digital twin architectures (e.g., Asset Administration Shell). I am eager to transition from purely generative physics to generative design assistance for secure systems.

\section{Proposed Research Directions}
Building upon the objectives of the **Security Engineering** project at CEA List/IRIT, I propose to investigate the following axes:

\subsection{1. AI-Assisted Security Modeling within MDE}
I intend to develop an AI-assisted environment that moves beyond simple prompting to become a systematic partner in the design process. By leveraging Large Language Models (LLMs) and GNNs, the framework will assist designers in identifying security vulnerabilities during the SysML/UML modeling phase and suggest standard-compliant mitigations in real-time.

\subsection{2. Automated Round-Trip Security Engineering}
The core of my research will focus on synchronizing design-time models with runtime execution feedback. I propose to develop mechanisms that automatically update architectural models based on security reports, vulnerability scans, and runtime attack detections. This "round-trip" capability ensures that the digital twin remains a living, secure representation of the physical infrastructure.

\subsection{3. Cyber-Security Simulation and Scenario Evaluation}
I propose to leverage the simulation capabilities of digital twins to evaluate complex attack scenarios. By injecting "virtual" threats into the digital twin's reference architecture, we can evaluate system configurations and mitigation strategies without risking the physical assets. This provides a safe "cyber-range" for testing the resilience of industrial infrastructures.

\section{Alignment with the EDT Program}
My background in AI-assisted modeling and security aligns with the goal of creating a secure national infrastructure for digital twin engineering. I am eager to contribute to the **Artemis platform** by providing a secure-by-design modeling layer that ensures the integrity and resilience of digital twins across the consortium.

\section{Conclusion}
The goal of my PhD will be to bridge the gap between AI-driven design assistance and rigorous security engineering. By formalizing an AI-assisted MDE framework for round-trip security, I aim to provide the foundational tools necessary for the next generation of resilient and self-securing industrial digital twins.

\end{document}
